Presenting users with raw numerical outputs, however, could result in significant misunderstanding.
To address this, we developed an LLM-based chatbot interface that supports natural-language question answering. 
When a user wants to know about a particular agent action, we construct a query prompt by collecting the current abstract state and the action taken by the agent. 
For counterfactual questions, we additionally extract the possible alternative actions from the agent policy.
We then instantiate one of two templates: ``Why does the agent choose [action] when in [state]?" (direct explanation) or ``Why does the agent choose [action] instead of choosing [alternate action] when in [state]?" (counterfactual explanation).
The LLM receives this prompt along with network information (e.g., machine connectivity) and the decision tree that encodes the alternative policy. 
It returns an explanation that describes why the action was taken or why an alternative was not taken in human-understandable natural language. The LLM does not generate explanations independently; instead, it translates structured outputs derived from SHAP and decision-tree policies into natural language.
In future work, for safety-critical deployment outside of the simulation environment, we plan to validate LLM outputs by ensuring consistency with the original feature set and prevent the introduction of unsupported information.

\textit{Example 3: Consider the query: ``Why does the blue agent remove access from 'user3' when there was a scan detected on 'enterprise1'?” The generated prompt is: ``Here is a decision tree showing the actions of a blue agent trying to stop a red agent on a network: [decision tree]. Here is information about the network: [network information]. Explain the following: ``Why does the agent remove 'user3' when 'enterprise1' is scanned and all other features are false?" The LLM responds: ``This could be for several reasons: (1). 'User3' might be communicating with 'Enterprise1', which could potentially be compromised by the red agent. By removing 'User3', the blue agent aims to cut off any potential communication between 'User3' and 'Enterprise1'. (2). 'User3' might be a vulnerable node in the network that could serve as an easy entry point for the red agent to move deeper into the network. By removing 'User3', the blue agent prevents this possibility. In summary, the blue agent removes 'User3' when there is no activity detected at 'Enterprise0', a scan is being performed on 'Enterprise1', and 'Enterprise1' shows some ongoing activity (indicating potential compromise or malicious action) as a precautionary measure to secure the network."}
